% Isolated proposed B171, checkpoint36. All old source bytes remain unchanged.
\subsection*{The exact limitation of the mirrored deletions}
\statement{B.171}{Proved}{The two entire mirrored one-sided fillings have identical restrictions on the actual fixed source, ordinarily and in its ambient compact closure, including every finite target refinement}
\label{stmt:B171}
In the actual B.123 filled group \(\Gamma_0=(F/N)\rtimes S\), put
\[
 \begin{aligned}
 L_{\mathrm{bar}}
   &=\langle\!\langle\bar a,\bar B,\bar d,\bar e
                    \rangle\!\rangle_{\Gamma_0},\\
 L_{\mathrm{plain}}
   &=\langle\!\langle a,B,d,e\rangle\!\rangle_{\Gamma_0}.
 \end{aligned}
\]
These are whole normal closures. Let
\(q_{\mathrm{bar}}\) and \(q_{\mathrm{plain}}\) be their quotient maps.
There is a marked quotient isomorphism
\begin{equation}\label{eq:B171-mirror-isomorphism}
 \theta:\Gamma_0/L_{\mathrm{plain}}
       \xrightarrow{\;\sim\;}\Gamma_0/L_{\mathrm{bar}},
 \qquad
 \theta q_{\mathrm{plain}}=q_{\mathrm{bar}}\delta.
\end{equation}
On the entire actual fixed subgroup, and hence on \(H_{\mathrm{src}}\),
the right side is \(q_{\mathrm{bar}}\). The same equality holds on the
actual compact fixed subgroup and actual source closure after completion.
In particular, both ordinary and compact restriction kernels on \(H_5\)
coincide. They are the normal defect groups identified in B.170.
Combining these two mirrored targets adds no source restriction.

\begin{proof}
The literal B.123.3 formulas, with \(\bar u=cuc\) and \(x=\tau c\), are
\[
 \delta(u)=x\bar ux^{-1},\qquad
 \delta(\bar u)=xux^{-1}
       \quad(u\in\{a,B,d,e,f\}),\qquad
 \delta(x)=x^{-1}.
\]
Both ordinary sheets of every node satisfy
\(\delta(r_0)=xr_1x^{-1}\) and
\(\delta(r_1)=xr_0x^{-1}\). Every signed source operator commutes with
\(\delta\). Thus the entire original normal closure \(N\) is preserved,
and these formulas define an involutive automorphism of the filled group.
They exchange \(L_{\mathrm{plain}}\) and \(L_{\mathrm{bar}}\), proving
\eqref{eq:B171-mirror-isomorphism}.

B.123.4 proves that the four actual contact words become fixed after
the node filling, while \(m,v\) are already fixed in the generic group.
Every element of \(H_{\mathrm{src}}\) is therefore fixed. On any fixed
element \(g\), equation \eqref{eq:B171-mirror-isomorphism} reads
\(\theta q_{\mathrm{plain}}(g)=q_{\mathrm{bar}}(g)\).
Since \(\theta\) is an isomorphism, the two kernels restricted to that
subgroup coincide.

Every map extends continuously to full profinite completions, and the
extension of \(\theta\) is again an isomorphism. The involution fixes
the actual source closure pointwise. The same calculation therefore
applies to actual compact closures, without identifying them with
completions of infinite-index subgroups. In particular it applies to
\(\mathcal H_5\), whose barred image has the full topology proved in
B.169, and B.170 identifies the common compact kernel as
\(\overline{\langle\!\langle[m,D_s]\rangle\!\rangle}_{\mathcal H_5}\).

For any finite target \(\varphi\) of the plain quotient, set
\(\varphi_{\mathrm{bar}}=\varphi\theta^{-1}\).
On the fixed source,
\[
 \varphi q_{\mathrm{plain}}
     =\varphi_{\mathrm{bar}}q_{\mathrm{bar}}.
\]
Conversely every finite barred target is obtained in this way.
The diagonal restriction into the product of the two quotient targets
is the graph of the same isomorphism, so also has the same kernel.
This proves the finite-refinement assertion as well as the kernel
equalities. It does not apply to unfilled contact words, which need
not be fixed before the disks are attached.

Thus the nontrivial defect \([m,D_s]\) is killed by both mirrored
deletions. A further comparison must retain information these targets
both erase. No failure of the actual full fixed-generation equality
is inferred.
\end{proof}
